\documentclass[aspectratio=169, 11pt]{beamer}

% --- Packages & Setup ---
\usepackage[utf8]{inputenc}
\usepackage[T1]{fontenc}
\usepackage{helvet}
\renewcommand{\familydefault}{\sfdefault} % Times New Roman font for text and math
\usepackage{xcolor}
\usepackage{booktabs}
\usepackage{graphicx}
\usepackage{tikz}
\usetikzlibrary{calc}

\graphicspath{{../}{./}}

% --- Branding Colors ---
\definecolor{amritaMaroon}{HTML}{A4123F}

% --- Theme Setup ---
\usetheme{Madrid}
\usecolortheme[named=amritaMaroon]{structure}
\setbeamertemplate{navigation symbols}{}
\setbeamertemplate{footline}[frame number]

\begin{document}

% =========================================================================
% TITLE SLIDE (Application-Based Project)
% =========================================================================
\begin{frame}
    % Corner Box Indicator
    \begin{tikzpicture}[remember picture,overlay]
        \node[fill=amritaMaroon, text=white, font=\bfseries\footnotesize, anchor=north east, rounded corners=2pt, inner sep=5pt]
        at ($(current page.north east) + (-0.2cm,-0.2cm)$)
        {Type 2: Application Project};
    \end{tikzpicture}

    \centering
    \includegraphics[width=2.8cm]{logo.png}
    \vspace{0.1cm}

    {\large \textbf{\textcolor{amritaMaroon}{[APPLICATION PROJECT TITLE]}}} \\
    \vspace{0.05cm}
    {\footnotesize \textbf{23CSE498 -- Project Phase II (External Review)} \quad | \quad \textbf{Team Number: [GXXX]}} \\
    \vspace{0.15cm}

    % Student Table
    \begin{table}
        \centering
        \footnotesize
        \renewcommand{\arraystretch}{1.05}
        \begin{tabular}{|c|c|l|}
            \hline
            \textbf{S.No} & \textbf{Register Number} & \textbf{Student Name} \\
            \hline
            1 & [CB.SC.U4CSE23XXX] & [Student Name 1] \\
            \hline
            2 & [CB.SC.U4CSE23XXX] & [Student Name 2] \\
            \hline
            3 & [CB.SC.U4CSE23XXX] & [Student Name 3] \\
            \hline
            4 & [CB.SC.U4CSE23XXX] & [Student Name 4] \\
            \hline
        \end{tabular}
    \end{table}
    \vspace{0.1cm}

    {\footnotesize Guide: \textbf{[Guide Name]}, [Designation], Department of CSE} \\
    \vspace{0.1cm}
    {\scriptsize \today}
\end{frame}

% =========================================================================
% SECTION 1: PROBLEM DEFINITION, REQUIREMENT ANALYSIS & SOLUTION DESIGN
% =========================================================================
% =========================================================================
% SLIDE 2: GUIDE APPROVAL EMAIL CONFIRMATION (Mandatory)
% =========================================================================
\begin{frame}{Guide Review \& Submission Approval}
    \framesubtitle{Confirmation of Project Review and Approval from Guide}
    \centering
    \fbox{\begin{minipage}[c][5.4cm]{0.94\textwidth}
        \centering
        \textbf{[Insert Screenshot of Guide Approval Email Here]} \\[0.3cm]
        {\footnotesize [Mandatory: Attach a clear, readable screenshot of the email received from your project guide approving this presentation and final review submission]}
    \end{minipage}}
\end{frame}

\section{Requirements \& Solution Design}

% --- SLIDE 1: PROBLEM DEFINITION & REQUIREMENT ANALYSIS ---
\begin{frame}{1. Problem Definition \& Requirement Analysis}
    \framesubtitle{Target Problem, User Needs \& Functional / Non-Functional Requirements}
    \begin{block}{Problem Statement \& Target Stakeholders}
        \begin{itemize}
            \item \textbf{Real-World Problem:} {[Define the practical problem or operational gap being solved]}
            \item \textbf{Target Users \& Use Context:} {[Identify primary end-users and application domain]}
        \end{itemize}
    \end{block}

    \begin{block}{Functional Requirements}
        \begin{itemize}
            \item \textbf{[FR-1]:} {[Specify primary user-facing functional requirement 1]}
            \item \textbf{[FR-2]:} {[Specify primary user-facing functional requirement 2]}
            \item \textbf{[FR-3]:} {[Specify primary user-facing functional requirement 3]}
        \end{itemize}
    \end{block}

    \begin{block}{Non-Functional Requirements \& Constraints}
        \begin{itemize}
            \item \textbf{[NFR-1]:} {[Specify responsiveness, usability, or reliability requirement]}
            \item \textbf{[NFR-2]:} {[Specify hardware, environmental, or operational constraint]}
        \end{itemize}
    \end{block}
\end{frame}

% --- SLIDE 2: END-TO-END APPLICATION ARCHITECTURE (FULL-SLIDE DIAGRAM) ---
% Note to students: Duplicate this frame if multiple architecture/workflow diagrams are required.
\begin{frame}{2. End-to-End Application Architecture}
    \framesubtitle{System Architecture \& User-to-Backend Workflow Diagram}
    \centering
    \fbox{\begin{minipage}[c][5.2cm]{0.94\textwidth}
        \centering
        \textbf{[Insert Full-Width Application Architecture / End-to-End Workflow Diagram]} \\[0.3cm]
        {\footnotesize [Placeholder: Illustrate User / Input Interface $\leftrightarrow$ Core Processing Modules $\leftrightarrow$ Data / Hardware / Output Interfaces]}
    \end{minipage}}
    \vspace{0.2cm}

    \begin{block}{Architecture Summary}
        \footnotesize {[Provide a concise 1--2 line walkthrough of how requests/inputs flow across the application components.]}
    \end{block}
\end{frame}

% --- SLIDE 3: ARCHITECTURAL COMPONENTS & INTERFACE DESIGN ---
% Note to students: Duplicate this frame if additional slides are needed to explain application sub-components.
\begin{frame}{3. Architectural Components \& Interface Design}
    \framesubtitle{Component Breakdown, Data Flow \& Technology Specifications}
    \begin{columns}[T]
        \column{0.50\textwidth}
        \begin{block}{Core Application Components}
            \footnotesize
            \begin{itemize}
                \item \textbf{[Component 1]:} {[Responsibility and functional scope]}
                \item \textbf{[Component 2]:} {[Responsibility and functional scope]}
                \item \textbf{[Component 3]:} {[Responsibility and functional scope]}
            \end{itemize}
        \end{block}

        \begin{block}{User Interaction \& Workflow Design}
            \footnotesize
            \begin{itemize}
                \item \textbf{Primary User Journey:} {[Outline step-by-step user interaction flow]}
                \item \textbf{Interface Design:} {[Describe key user views or control interfaces]}
            \end{itemize}
        \end{block}

        \column{0.46\textwidth}
        \begin{block}{Data Management \& External Interfaces}
            \footnotesize
            \begin{itemize}
                \item \textbf{Data Storage:} {[Structure of stored records / state management]}
                \item \textbf{External Interfaces:} {[Hardware, sensors, or external service interfaces]}
            \end{itemize}
        \end{block}

        \begin{block}{Platform \& Deployment Specifications}
            \footnotesize
            \begin{itemize}
                \item \textbf{Development Stack:} {[Selected tools and frameworks]}
                \item \textbf{Deployment Target:} {[Target execution platform]}
            \end{itemize}
        \end{block}
    \end{columns}
\end{frame}

% =========================================================================
% SECTION 2: IMPLEMENTATION COMPLETION (100%)
% =========================================================================
\section{Implementation Details}

% --- SLIDE 4: IMPLEMENTATION COMPLETION (100%) - OVERVIEW ---
\begin{frame}{4. Implementation Details -- Overview}
    \framesubtitle{End-to-End Application Module \& Integration Status}
    \begin{table}
        \centering
        \footnotesize
        \renewcommand{\arraystretch}{1.25}
        \begin{tabular}{p{3.6cm} p{6.2cm} c}
            \toprule
            \textbf{Application Module} & \textbf{Implemented Features \& Integration Scope} & \textbf{Status} \\
            \midrule
            {[Module 1: User Interface]} & {[Summary of completed user views, forms, and interactive controls]} & 100\% (Integrated) \\
            {[Module 2: Core Domain Logic]} & {[Summary of implemented application workflows and processing logic]} & 100\% (Integrated) \\
            {[Module 3: Data \& Persistence]} & {[Summary of implemented storage, retrieval, and state handling]} & 100\% (Integrated) \\
            {[Module 4: External Integration]} & {[Summary of integrated external services, devices, or notifications]} & 100\% (Functional) \\
            \bottomrule
        \end{tabular}
    \end{table}

    \begin{block}{Prototype Environment \& Repository}
        \footnotesize
        \begin{itemize}
            \item \textbf{Execution Platform:} {[Specify target runtime environment where the complete application is deployed/executed]}
            \item \textbf{Source Repository:} {[Insert repository URL / commit reference for the integrated application codebase]}
        \end{itemize}
    \end{block}
\end{frame}

% --- SLIDE 5: FUNCTIONAL MODULE IMPLEMENTATION DETAILS ---
% Note to students: Duplicate this frame if needed to explain additional implemented application modules.
\begin{frame}{5. Functional Module Implementation Details}
    \framesubtitle{Key Feature Implementation \& Inter-Module Integration}
    \begin{columns}[T]
        \column{0.50\textwidth}
        \begin{block}{[Functional Module A Name]}
            \footnotesize
            \begin{itemize}
                \item \textbf{Implemented Workflow:} {[Explain how this feature works end-to-end]}
                \item \textbf{Inputs \& Outputs:} {[Specify user/system inputs and generated outputs]}
                \item \textbf{Validation \& Controls:} {[Describe input checks and state updates]}
            \end{itemize}
        \end{block}

        \begin{block}{[Functional Module B Name]}
            \footnotesize
            \begin{itemize}
                \item \textbf{Implemented Workflow:} {[Explain how this feature works end-to-end]}
                \item \textbf{Inter-Module Integration:} {[Describe how it connects with other modules]}
            \end{itemize}
        \end{block}

        \column{0.46\textwidth}
        \centering
        \vspace{0.2cm}
        \fbox{\begin{minipage}[c][5.0cm]{0.94\textwidth}
            \centering
            \textbf{[Insert Module Integration Flow / Feature Implementation Snapshot]} \\[0.2cm]
            {\scriptsize [Placeholder: Show connected module execution or key functional view]}
        \end{minipage}}
    \end{columns}
\end{frame}

% =========================================================================
% SECTION 3: FUNCTIONAL TESTING & DEMONSTRATION
% =========================================================================
\section{Functional Testing \& Demonstration}

% --- SLIDE 6: FUNCTIONAL TESTING & VERIFICATION RESULTS ---
\begin{frame}{6. Functional Testing \& Verification Results}
    \framesubtitle{End-to-End Feature Verification \& Test Case Summary}
    \begin{table}
        \centering
        \footnotesize
        \renewcommand{\arraystretch}{1.25}
        \begin{tabular}{p{3.4cm} p{3.2cm} p{3.6cm} c}
            \toprule
            \textbf{Test Case / Feature} & \textbf{Test Input / Action} & \textbf{Verified Outcome} & \textbf{Status} \\
            \midrule
            {[TC-01: Primary Workflow]} & {[Input / User Action]} & {[Expected output verified]} & {[Pass]} \\
            {[TC-02: Secondary Workflow]} & {[Input / User Action]} & {[Expected output verified]} & {[Pass]} \\
            {[TC-03: Data / Device Sync]} & {[Input / User Action]} & {[Expected output verified]} & {[Pass]} \\
            {[TC-04: Boundary / Error Case]} & {[Invalid / Edge Input]} & {[Graceful handling verified]} & {[Pass]} \\
            \bottomrule
        \end{tabular}
    \end{table}

    \begin{block}{Application Verification Summary}
        \footnotesize
        \begin{itemize}
            \item \textbf{Functional Verification:} {[Summarize end-to-end feature verification across target use cases]}
            \item \textbf{Responsiveness \& Operational Metrics:} {[Summarize latency, accuracy, or user testing observations]}
        \end{itemize}
    \end{block}
\end{frame}

% --- SLIDE 7: WORKING PROTOTYPE DEMONSTRATION ---
% Note to students: Use a single full-width box or two side-by-side boxes as appropriate for your prototype screenshots.
\begin{frame}{7. Working Prototype Demonstration}
    \framesubtitle{End-to-End Application Walkthrough \& Visual Evidence}
    \begin{columns}[T]
        \column{0.48\textwidth}
        \centering
        \fbox{\begin{minipage}[c][4.8cm]{0.94\textwidth}
            \centering
            \textbf{[Insert Working Prototype Screenshot 1: Input / Control / Setup View]} \\[0.2cm]
            {\scriptsize [Placeholder: Primary user interface or hardware prototype setup]}
        \end{minipage}}

        \column{0.48\textwidth}
        \centering
        \fbox{\begin{minipage}[c][4.8cm]{0.94\textwidth}
            \centering
            \textbf{[Insert Working Prototype Screenshot 2: Processed Output / Action View]} \\[0.2cm]
            {\scriptsize [Placeholder: End-to-end result screen or physical output demonstration]}
        \end{minipage}}
    \end{columns}
    \vspace{0.2cm}

    \begin{block}{Live Demonstration Scenario}
        \footnotesize {[Briefly describe the end-to-end user scenario demonstrated live during the review.]}
    \end{block}
\end{frame}

% =========================================================================
% SECTION 4: TECHNICAL KNOWLEDGE & VIVA VOCE
% =========================================================================
\section{Technical Knowledge}

% --- SLIDE 8: TECHNICAL JUSTIFICATION & ARCHITECTURAL DECISIONS ---
\begin{frame}{8. Technical Justification \& Architectural Decisions}
    \framesubtitle{Platform Rationale, Engineering Trade-offs \& Operational Considerations}
    \begin{columns}[T]
        \column{0.48\textwidth}
        \begin{block}{1. Platform \& Technology Justification}
            \footnotesize
            \begin{itemize}
                \item {[Explain why the chosen tools, frameworks, or hardware components best suit the application requirements]}
                \item {[Compare with alternative choices considered]}
            \end{itemize}
        \end{block}

        \begin{block}{2. Architecture \& Integration Trade-offs}
            \footnotesize
            \begin{itemize}
                \item {[Discuss design choices balancing responsiveness, resource usage, and modularity]}
                \item {[Explain inter-module communication reliability]}
            \end{itemize}
        \end{block}

        \column{0.48\textwidth}
        \begin{block}{3. Usability, Security \& Safety Measures}
            \footnotesize
            \begin{itemize}
                \item {[Detail user data protection, access control, and input validation mechanisms]}
                \item {[Describe error recovery and fallback handling]}
            \end{itemize}
        \end{block}

        \begin{block}{4. Operational Constraints \& Future Scope}
            \footnotesize
            \begin{itemize}
                \item {[State practical operating boundaries of the current prototype]}
                \item {[Outline planned enhancements for deployment scale]}
            \end{itemize}
        \end{block}
    \end{columns}
\end{frame}

% =========================================================================
% REFERENCES & THANK YOU
% =========================================================================
\section{References}
\nocite{*}
\begin{frame}[allowframebreaks]{References}
    \bibliographystyle{ieeetr}
    \bibliography{ref}
\end{frame}

\begin{frame}
    \centering
    {\Huge \textbf{\textcolor{amritaMaroon}{Thank You}}} \\[0.4cm]
    {\large \textbf{[Questions \& Viva Voce]}}
\end{frame}

\end{document}
